%% kapa: exam=1 round=34 year=2023 period=2 form=A subject=accounting subject_ko=회계학 seq=23 number=63
%% source: Q-Net 감정평가사 기출문제 (cst00309) artlSeq=5211789 file=raw/1008/5211789/2023년도 제34회 감정평가사 1차시험 2교시.pdf
%% conversion: kordoc 4.24.51 -> tools/gen1.py (auto); review=proofread-sonnet; flags=none
\begin{question}{63}
고객과의 계약에서 생기는 수익에 관한 설명으로 옳은 것은?
\begin{choices}
  \item 계약의 결과로 기업의 미래 현금흐름의 위험, 시기, 금액이 변동될 것으로 예상되지 않는 경우에도 고객과의 계약으로 회계처리할 수 있다.
  \item 계약은 서면으로, 구두로, 기업의 사업 관행에 따라 암묵적으로 체결할 수 있다.
  \item 이전할 재화나 용역의 지급조건을 식별할 수 없는 경우라도 고객과의 계약으로 회계처리할 수 있다.
  \item 계약변경은 반드시 서면으로만 승인될 수 있다.
  \item 고객과의 계약에서 식별되는 수행의무는 계약에 분명히 기재한 재화나 용역에만 한정된다.
\end{choices}
\end{question}
\answer{2}
